\scanpage{167}
\mspage{80}
avec cette notation, l'\uline{âme} du topos \inserted{élémentaire} envisagé est $D^*_{I\setminus I'_1}$, qu'on va considérer comme plongé dans $D^*_{I'_0}$
\[
(271)\qquad
\underset{\text{(``âme'')}}{A}=D^*_{\underbrace{I\setminus I'_1}_{I'_0\amalg I''}}
\underset{\substack{\text{immersion}\\\text{de codim }I''}}{\subset}
D^*_{I'_0}=\underset{\text{(``réceptacle'')}}{R}
\]
Le \deleted{\unsure{premier}} cas où on \unsure{aurait|a} $I''=I''_0$ correspond à celui du voisinage tubulaire de \deleted{$\unsure{D^*_{I\setminus I'_1}}$ dans $D_{I'_0}$} $A$ dans $R$, soit $\unsure{\mathcal V_{A,R}|\nu_{A,R}}$, celui où $I''=\unsure{I''_*|I''{}^*}$ correspond au cas de $\unsure{\mathcal V^*_{A,R}|\nu^*_{A,R}}$ (le $*$ bien sûr relatif ici à la situation induite sur $R$, ou si on préfère, sur $D_{I'_0}\supset R$, avec le diviseur induit par $\Theta_{d_0}$ \ldots). Le cas intermédiaire, ``général'', correspond donc à \deleted{celui} la donnée d'une strate intermédiaire entre $A$ et $R$
\[
(272)
\]
\diagram{\begin{tikzcd}[ampersand replacement=\&,column sep=small,row sep=large]
A \arrow[rr,phantom,"\subset"] \arrow[d,equal]
\&\&B \arrow[rr,phantom,"\subset"] \arrow[d,equal]
\&\&R \arrow[d,equal]\\
D^*_{I'_0\amalg I''}\&{}\arrow[d,no head]
\&D^*_{I'_0\amalg I''{}^*}\&{}\arrow[d,no head]\&D_{I'_0}\\
\&\text{codim }\operatorname{card}I''_0=d''_0
\&\&\text{codim}=\operatorname{card}I''{}^*=d''{}^*\&
\end{tikzcd}}
\margin{166}
